\begin{frame}{JSON, Javascript Object Notation}

\begin{itemize}
\item JavaScript Object Notation;
\item Initialement créé pour la sérialisation et l'échange d'objets JavaScript;
\item Langage pour l'échange de données semi-structurées (et éventuellement structurées);
\item Format texte indépendant du language de programmation utilisé pour le manipuler.
\end{itemize}

\vfill 

Utilisation première : échange de données dans un environnement Web (par exemple applications Ajax)\\

\vfill

Extension : sérialisation\prof{ pour échange évidemment} et stockage de données \prof{semi-structurées plut\^ot mais aussi structurées si besoin}

\end{frame}

\definecolor{lightgray}{rgb}{.9,.9,.9}
\definecolor{darkgray}{rgb}{.4,.4,.4}
\definecolor{purple}{rgb}{0.65, 0.12, 0.82}


\begin{frame}[containsverbatim,fragile,allowframebreaks]{Les bases de JSON}

Structure de base: paire \prof{ordonnée} \textit{clef-valeur} (\textit{key-value})


\begin{lstlisting}[language=JavaScript,basicstyle=\small]
"title": "The Social network"
\end{lstlisting}

\bigskip

\alert{Qu'est-ce qu'une valeur?} On distingue les
 valeurs \alert{atomiques} et les valeurs \alert{complexes} (construites)

\medskip
Valeurs atomiques\,: cha\^ines de caractères (entourées 
par les classiques guillemets anglais (droits)), nombres (entiers, flottants) 
et valeurs booléennes (\texttt{true} ou \texttt{false}).

\begin{lstlisting}[language=JavaScript,basicstyle=\small]
"year": 2010
\end{lstlisting}


\begin{lstlisting}[language=JavaScript,basicstyle=\small]
"oscar": false
\end{lstlisting}

\framebreak

Valeurs complexes\,: les \alert{objets}.

\bigskip

Un \textit{objet} est un ensemble de paires clef-valeur.

\bigskip

Au sein d'un ensemble de paires, une clef apparait au plus une fois 
(NB : les types de valeurs peuvent être distincts).

\begin{lstlisting}[language=JavaScript,basicstyle=\small]
{"last_name": "Fincher", "first_name": "David"}
\end{lstlisting}

\bigskip

Un objet peut être utilisé comme valeur (dite \textit{complexe}) dans une paire clef-valeur.

\begin{lstlisting}[language=JavaScript,basicstyle=\small]
"director": { 
  "last_name": "Fincher",
  "first_name": "David", 
  "birth_date": 1962
}
\end{lstlisting}

\framebreak

Valeurs complexes\,: les \alert{tableaux}.

\bigskip

Un \textit{tableau} (\textit{array}) est une liste 
 de valeurs (dont le type n'est pas forcément le même).

\begin{lstlisting}[language=JavaScript,basicstyle=\small]
"actors": ["Eisenberg", "Mara", "Garfield", "Timberlake"]
\end{lstlisting}

\bigskip

Imbrication sans limite (vous vous souvenez d'XML?): tableaux de tableaux, tableaux d'objets contenant eux-mêmes des tableaux, etc.
\framebreak

Un \textit{document} est un objet. Il peut être défini par des objets et tableaux imbriqu\' es autant de fois que nécessaire.

\begin{lstlisting}[language=JavaScript,basicstyle=\small]
{
  "title": "The Social network", 
  "summary": "On a fall night in 2003, Harvard undergrad and \n
         programming genius Mark Zuckerberg sits down at his \n 
         computer and heatedly begins working on a new idea. (...)",
  "year": 2010, 
  "director": {"last_name": "Fincher",
                    "first_name": "David"},
  "actors": [ 
        {"first_name": "Jesse", "last_name": "Eisenberg"}, 
        {"first_name": "Rooney", "last_name": "Mara"}
    ]
}
\end{lstlisting}

\end{frame}


\begin{frame}{JSON vs. XML}
\begin{itemize}
\item JSON plus léger et intuitif que XML,
\item Facile à parser pour n'importe quel langage de programmation,
\item JSON n'a pas (encore) de langage de spécification de schéma associé,
\item JSON n'a pas (encore) de langage de requête associé. \prof{moins mature en quelque sortes}
\end{itemize}
\end{frame}



\begin{frame}{Quelques sites pour aler plus loin}
\begin{itemize}
\item Tout sur JSON\,: \url{http://json.org/}
\item Un validateur de documents JSON\,: \url{http://jsonlint.com/}
\item Une proposition de \alert{schéma} pour JSON\,: \url{http://json-schema.org/}
\item Un langage de requêtes  JSON\,: JAQL, \url{http://code.google.com/p/jaql/}


\end{itemize}
\end{frame}


